\chapter{Get familiar with the project}

    \section{KEcoLab repository}
    \textbf{What's inside of the repository? (structure, actual code, etc.)}

    \section{ydotool, kdotool and xdotool - Actual problems}
    \textbf{How does ydotool work? Pros, Cons?}

    \section{The scripting scenarios}
    \textbf{Description of the scenarios, what are the steps, which are done}

    \subsection{Baseline scenario}
    \subsection{Idle Scenario}
    \subsection{System Usage Scenario}

    \section{Setting up the environment}
    \textbf{What is the working setup generally?}

    \subsection{Fedora VM with KDE Plasma}
    % gobble here takes away n amount of space from each code block line
    \begin{lstlisting}[gobble=8] 
        sudo dnf upgrade
    \end{lstlisting}

    \subsubsection{Setup ydotool and kdotool}
    To run the actual scripts, also kdotool and ydotool have to be set up, as 
    the scripts are built on these programs. So first for the installation 
    make sure, to also have cmake, g++/clang and sdoc installed. If not, 
    install them(in this example g++ is used, to install clang just swap g++ 
    with clang):
    \begin{lstlisting}[gobble=8]
        sudo dnf install cmake g++ scdoc
    \end{lstlisting}

    Further clone the repository and install the application by using:
    \begin{lstlisting}[gobble=8]
        git clone -b feat/setup-script https://invent.kde.org/neogg/ydotool.git
        cd ydotool
        chmod +x setup.sh
        ./setup.sh install
    \end{lstlisting}

    Also give ydotoold persmissions to uinput:
    \begin{lstlisting}[gobble=8]
        sudo touch /etc/udev/rules.d/99-uinput.rules
        echo 'KERNEL=="uinput", MODE="0660", GROUP="input"' | sudo tee /etc/udev/rules.d/99-uinput.rules > /dev/null
        sudo usermod -aG input "$USER"

        sudo udevadm control --reload-rules
        sudo udevadm trigger
    \end{lstlisting}

    Now the system has to be rebooted and after that ydotool can be started by 
    using:
    \begin{lstlisting}[gobble=8]
        systemctl --user start ydotoold
    \end{lstlisting}

    To verify, that everything worked you can try some ydotool commands, 
    as example lets try to do the bash command ls via the  ydotool interface:
    \begin{lstlisting}[gobble=8]
        ydotool type "ls" && ydotool key 28:1 28:0
    \end{lstlisting}
    This types "ls" to the command-line and key does press keys - 
    28 stands for enter, 0 is release and 1 is press.

    A keymap where all keys are described can be found in `/usr/include/linux/input-event-codes.h`.

    \textbf{!!!MISSING KDOTOOL!!!}
    \newpage

    \subsection{Run the project locally and test the scripts}
    Now that the VM is setup and ydotool/kdotool are installed, it can be tried to run the actual scripts. As the script of Okular is the only script, which is working, those will be run. First, the KEcoLab repository has to be cloned locally by using:
    \begin{lstlisting}[gobble=8]
        git clone https://invent.kde.org/sdk/kecolab.git
    \end{lstlisting}
    This will create the project folder, in order to be able to access the scripts. These are stored in `scripts/test\textunderscore scripts/org.kde.okular`. The folder contains three scripting files, one for each scenario.
    To be able to run the scripts you need the flatpak version of the application. As this script is now for Okular, the `org.kde.okular` package
    needs to be installed by running `flatpak install org.kde.okular`.
    If not installed, the script is not able to start the application.
    Also the scripts have to be executable, which is done by `chmod u+x $<$file-name$>$`. After these steps the script can be started and will run.
    For example starting the idle script would be `./log\textunderscore idle.sh`

    The scripts have principially the same structure until it comes to the testing for-loop. Each one has a syncUp-function defined, which sets sleep for the given time.
    `startAction()` and `stopAction()` will log the exact time of each action's start and stop to `~/log\textunderscore usage\textunderscore scenario`, as each action is wrapped with both functions. These functions are only defined in sus and idle scenario, as with baseline-scenario it's not needed.

    \subsubsection{Baseline scenario}
    This scenario is the most simple one, as there's not application started or anything else done, than logging the time of each iteration's start/end and waiting for 218s (including burn-in time of 60s and cooldown of 30s)

    \subsubsection{Idle Scenario}
    The idle scenario additionally contains the start of the application. Furthermore it will let Okular run for 220 seconds and then close it, which makes the scenario of the plain application running, with no actions done.

    \subsubsection{System usage Scenario}
    For running the system usage scenario, the configuration of the following key-shortcuts have to be made, as ydotool does many operations using shortkeys:
    \begin{itemize}
        \item Rotate Left: Ctrl+l
        \item Rotate Right: Ctrl+r
        \item Invert color: Ctrl+i
        \item Fit to width: Ctrl+Shift+w
        \item Single page mode: Ctrl+Alt+s
    \end{itemize}
    Also make shure that ydotool is running.
    The lines until the main for-loop are basically the same as in the idle script. For simulating the usage, a PDF-file is needed. The existance of a testing file is checked during the script at the beginning, and downloaded, if the file is not at the specified path. As the for loop starts, the pdf gets copied to the home directory:
    \begin{lstlisting}[gobble=8]
        cp ~/Documents/okular/20yearsofKDE.pdf ~/Documents/20yearsofKDE
    \end{lstlisting}
    After copying the file and a wait time of 60s the file gets opened with:
    \begin{lstlisting}[gobble=8]
        flatpak run org.kde.okular ~/Documents/20yearsofKDE > /dev/null 2>&1 &
    \end{lstlisting}
    `/dev/null 2$>$\&1 \&`suspends the output, so no output is given to the command line.

    \subsection{Change the source code and perform a merge request}